% Figura: e^{-1/x^2}, función C^infinito no analítica
\begin{figure}[ht]
\centering
\begin{tikzpicture}
  \begin{groupplot}[group style={group size=2 by 1, horizontal sep=1.6cm},
      oscuro, width=0.47\textwidth, height=5.0cm, samples=301]
    \nextgroupplot[xmin=-4, xmax=4, ymin=-0.3, ymax=1.08,
        xtick={-4,-2,0,2,4}, ytick={0,0.5,1},
        xlabel={$x$}, ylabel={$f(x)$},
        title={(a) Vista global}]
      \addplot[c1, domain=-4:4] {exp(-1/(x^2+1e-12))};
      \addplot[c3, densely dashed, domain=-4:4] {0};
      \addplot[guia, domain=-4:4, forget plot] {1};
      \node[font=\scriptsize, text=cian, anchor=west] at (axis cs:2.05,0.55) {$\eul^{-1/x^{2}}$};
      \node[font=\scriptsize, text=ambar, anchor=north] at (axis cs:0,-0.06)
        {serie de Taylor en $0$: $\equiv0$};
    \nextgroupplot[xmin=-0.5, xmax=0.5, ymin=-0.0055, ymax=0.02,
        xtick={-0.5,-0.25,0,0.25,0.5},
        ytick={0,0.01,0.02},
        yticklabel style={/pgf/number format/fixed, /pgf/number format/precision=2},
        scaled y ticks=false,
        xlabel={$x$},
        title={(b) Detalle cerca del origen}]
      \addplot[c1, domain=-0.5:0.5] {exp(-1/(x^2+1e-12))};
      \addplot[c3, densely dashed, domain=-0.5:0.5] {0};
  \end{groupplot}
\end{tikzpicture}
\caption{La función $f(x)=\eul^{-1/x^{2}}$ ($f(0)=0$) es de clase
$\Cinf$ y todas sus derivadas se anulan en el origen, de modo que su serie de
Taylor en $0$ es la función nula (trazo discontinuo). En (b) se aprecia lo
extraordinariamente plana que es $f$ cerca de $0$: en $x=\pm0{,}25$ vale
apenas $\eul^{-16}\approx1{,}1\cdot10^{-7}$.}
\label{fig:suave}
\end{figure}
